% K31-30, JC-003 y P01 del inventario REV05.
% Residencia: después de la derivación del registro K y de la incidencia Witt;
% antes de utilizar la lectura ponderada para la recuperación del registro.
% Propietarios: output/IMPLEMENTACION_K_SEMILLAS_INCIDENCIA_20260918/
% WeightedIncidenceRecovery.lean (íntegro) y
% COMPOSICION_GENERATIVA_Y_RECUPERACION_K.tex:257--306.
% Codificador: PAQUETE_ARTICULO_I_INTEGRACION_ESTADO_CAMPOS_20260919/
% lean/biblioteca/CoxeterNeighbor.lean:322--343;
% lean/biblioteca/SectorIncidenceData.lean:27--33.
% Estatuto: RESULTADO_RECUPERADO, explicitación íntegra de las matrices que
% la exposición anterior remitía al código. La prueba Lean es la ya existente.

\section[Recovery through hexadic incidences]{Complete recovery of the register through weighted hexadic incidences}
\label{rh:sec:recuperacion}

The generated dodecaphasic register preserves twelve ordered components and
a phase origin. Paley--Witt incidence provides an additional
reading of that register: each hexadic support combines six of its
components through a sum weighted by the incidence itself. The
reconstruction developed here receives the values of these sums
from the same register. Its purpose is to prove that the readings preserve
all the components and to exhibit the inverse operator with explicit
rational coefficients.

\subsection{Incidence encoder and twelve explicit witnesses}

In the incidence coordinate system already fixed, let
\[
A_W=\begin{pmatrix}
0&1&1&1&1&1\\1&0&1&2&2&1\\1&1&0&1&2&2\\
1&2&1&0&1&2\\1&2&2&1&0&1\\1&1&2&2&1&0
\end{pmatrix}\in M_6(\mathbb F_3).
\]
For $w\in\mathbb F_3^6$, the encoder is
$\operatorname{Enc}(w)=(w,wA_W)\in\mathbb F_3^{12}$ and its support
is
\[
H(w)=\{j\in\{1,\ldots,12\}:\operatorname{Enc}(w)_j\ne0\}.
\]
Choosing a position uses indices from one to twelve;
implementations with zero-based indices are transported through
$j\mapsto j+1$.

The following twelve witnesses belong to the domain of the encoder:
\[
\begin{array}{c|c@{\qquad}c|c}
i&w^{(i)}&i&w^{(i)}\\\hline
1&112210&7&122020\\
2&122101&8&112002\\
3&121200&9&111001\\
4&111100&10&110122\\
5&121012&11&101221\\
6&111020&12&011111.
\end{array}
\]
Define $H_i=H(w^{(i)})$ and
$B_{ij}=\mathbf1_{\{j\in H_i\}}$. Evaluating the encoder gives
\begingroup\small\setlength{\arraycolsep}{3pt}
\setcounter{MaxMatrixCols}{12}
\begin{equation}
B=\begin{pmatrix}
1&1&1&1&1&0&0&0&0&0&0&1\\
1&1&1&1&0&1&0&0&0&0&1&0\\
1&1&1&1&0&0&1&0&1&0&0&0\\
1&1&1&1&0&0&0&1&0&1&0&0\\
1&1&1&0&1&1&0&0&0&1&0&0\\
1&1&1&0&1&0&1&0&0&0&1&0\\
1&1&1&0&1&0&0&1&1&0&0&0\\
1&1&1&0&0&1&1&1&0&0&0&0\\
1&1&1&0&0&1&0&0&1&0&0&1\\
1&1&0&1&1&1&0&0&1&0&0&0\\
1&0&1&1&1&1&0&1&0&0&0&0\\
0&1&1&1&1&1&1&0&0&0&0&0
\end{pmatrix}.
\label{rh:eq:incidencia}
\end{equation}
\endgroup
Each row contains six entries equal to one. Consequently, each
$H_i$ is a hexad of the image of the encoder, and the matrix $B$ is
constructed from those supports.

\subsection{Rational inverse and uniqueness of recovery}

Let the integer matrix be
\begingroup\small\setlength{\arraycolsep}{3pt}
\setcounter{MaxMatrixCols}{12}
\begin{equation}
M=\begin{pmatrix}
7&1&8&-7&7&-1&-4&8&-7&3&3&-15\\
7&7&-4&-1&1&-7&8&8&-7&3&-15&3\\
1&7&8&-7&7&-7&8&-4&-1&-15&3&3\\
1&1&2&5&-5&-1&-4&-4&-1&3&3&3\\
1&-5&-4&-1&1&5&2&-4&-1&3&3&3\\
-5&1&-4&-1&1&-1&-4&2&5&3&3&3\\
-5&-11&2&5&-5&11&-10&2&5&3&3&3\\
-5&-5&-10&11&-11&5&2&2&5&3&3&3\\
-11&-5&2&5&-5&5&2&-10&11&3&3&3\\
-11&-11&-4&17&1&11&-10&-10&11&3&3&3\\
-11&1&-10&11&-11&17&-4&-10&11&3&3&3\\
1&-11&-10&11&-11&11&-10&-4&17&3&3&3
\end{pmatrix}.
\label{rh:eq:inversa-entera}
\end{equation}
\endgroup

\begin{theorem}[Exact recovery from twelve hexadic readings]
\label{rh:thm:doce}
For $v\in\mathbb Q^{12}$, let
\[
\operatorname{Read}(v)=Bv,
\qquad
\operatorname{Rec}(y)=\frac1{18}My.
\]
We verify
\[
\operatorname{Rec}(\operatorname{Read}(v))=v,
\qquad
\operatorname{Read}(\operatorname{Rec}(y))=y.
\]
Therefore every $y\in\mathbb Q^{12}$ admits a unique register
$v\in\mathbb Q^{12}$ with $\operatorname{Read}(v)=y$.
\end{theorem}
\begin{proof}
Integer multiplication of the matrices
\eqref{rh:eq:incidencia} and \eqref{rh:eq:inversa-entera} gives
$MB=BM=18I_{12}$. Each diagonal entry results from a scalar
product equal to eighteen and each off-diagonal entry results
from a zero scalar product. Dividing by eighteen gives
the two inverse identities.

In particular,
$\operatorname{Rec}(\operatorname{Read}(v))=(MB/18)v=v$ and
$\operatorname{Read}(\operatorname{Rec}(y))=(BM/18)y=y$.
The second identity proves existence through
$v=\operatorname{Rec}(y)$. If $Bu=Bv$, applying $M/18$ to both
sides gives $u=v$, which proves uniqueness.
\end{proof}

The check of $B$ from the encoder and the two exact
multiplications are formalized in the theorems
\texttt{incidence\_from\_encoder}, \texttt{inverse\_left} and
\texttt{inverse\_right}. The composition identities correspond
to \texttt{recover\_read} and \texttt{read\_recover}.
Their mathematical content is found entirely in the matrices and
in the preceding proof; the formal certificate records the verification
of those same operations.

\begin{corollary}[Separation by all weighted incidences]
If $u,v\in\mathbb Q^{12}$ satisfy
\[
\sum_{j\in H(w)}u_j=\sum_{j\in H(w)}v_j
\]
for every $w\in\mathbb F_3^6$ whose encoded support has six
positions, then $u=v$.
\end{corollary}
\begin{proof}
The hypothesis applies in particular to the twelve displayed witnesses.
By the definition of $B$, their twelve equalities constitute $Bu=Bv$.
Theorem~\ref{rh:thm:doce} allows both inputs to be recovered and
concludes $u=v$.
\end{proof}

\begin{example}[Recovery of a peripheral component]
For the canonical vector $v=e_{12}$, the only nonzero readings are
the first and the ninth:
\[
y=Be_{12}=(1,0,0,0,0,0,0,0,1,0,0,0)^{\mathsf T}.
\]
The sum of the first and ninth columns of $M$ is $18e_{12}$.
Consequently, $My/18=e_{12}$. The two hexadic sums distinguish
this position in the joint reading because the remaining components
cancel exactly in the recovery operator.
\end{example}

\subsection{Symmetric recovery operator on the one hundred and thirty-two hexads}

The Steiner system $S(5,6,12)$ already produced by incidence
contains $132$ hexads. Denote their set by $\mathcal H$ and
define
\[
\mathcal I:\mathbb Q^{12}\longrightarrow\mathbb Q^{\mathcal H},
\qquad
(\mathcal I v)_H=\sum_{j\in H}v_j.
\]
This reader has redundancy: it preserves the twelve components through
$132$ sums, whereas the preceding theorem provides a basis of
twelve of these sums.

\begin{theorem}[Inverse on the image of the complete reader]
\label{rh:thm:steiner}
We have
\[
\mathcal I^{\mathsf T}\mathcal I=36I_{12}+30J_{12},
\]
where $J_{12}$ has all its entries equal to one. For
$y=\mathcal I v$, the register is recovered through
\begin{equation}
396v_i=11\sum_{H\ni i}y_H-5\sum_{H\in\mathcal H}y_H,
\qquad i=1,\ldots,12.
\label{rh:eq:recuperador}
\end{equation}
\end{theorem}
\begin{proof}
The property $S(5,6,12)$ implies that each position belongs to
$\binom{11}{4}/\binom{5}{4}=66$ hexads. Each pair of distinct
positions belongs to $\binom{10}{3}/\binom{4}{3}=30$ hexads.
Therefore, the product of the transposed incidence matrix with
itself has diagonal $66$ and off-diagonal entries $30$, which
is $36I+30J$.

Let $S=\sum_jv_j$. The two preceding counts give
\[
\sum_{H\ni i}y_H=66v_i+30\sum_{j\ne i}v_j=36v_i+30S,
\qquad
\sum_Hy_H=66S.
\]
Multiplying the first equality by eleven and subtracting five times the
second eliminates $S$ and produces $396v_i$. This proves
\eqref{rh:eq:recuperador} for every reading belonging to the
image of $\mathcal I$.
\end{proof}

\begin{corollary}[Uniform and transverse components]
On $\mathbb R^{12}$, the complete reader has singular value
$\sqrt{396}$ in the uniform direction and singular value $6$ in its
orthogonal complement. Its inverse operator on the image has operator
norm $1/6$.
\end{corollary}
\begin{proof}
The matrix $J$ acts by twelve on $\mathbb R\mathbf1$ and by zero
on $\mathbf1^\perp$. The eigenvalues of $36I+30J$ are therefore
$396$ and $36$. The singular values are their positive
roots; the largest singular value of the inverse operator is the reciprocal
of the smallest of them.
\end{proof}

\subsection{Composition with the generated register}

Applying these identities to the integer dodecaphasic register $K$, already
produced with its order and origin, the readings $BK$ or $\mathcal I K$
allow its twelve components to be recovered exactly. The subsequent positional
interpretation takes that same register in base one thousand; the
Hadamard operators, differences and hexadic readers are
recoverable realizations of a common provenance.

The system of twelve rational coordinates admits any input
$y\in\mathbb Q^{12}$ and reconstructs it as a unique rational
register. The subdomain corresponding to generated HMT registers
additionally preserves its conditions of integrality, triad
width, origin, selection and boundary. For the redundant system of
$132$ readings, the inverse acts on the image of $\mathcal I$.
These specifications type the recovered content without modifying the
antecedent generation of the register.
